Read in Gauge Theory Gravity, the geometric steps below (Ricci curvature, the local boost field, the field equations) are those of the gravitational field on flat spacetime through its derived metric and field strength (Fact \claimref{synthesis.imported-gtg-dictionary}, clauses (iv) and (vi)). \emph{(i)} Let $\mathfrak{A}_{F}$ be a window of the modes of a local horizon through a point of $U$, with the boost acting by $a_{j} \mapsto e^{-Jsw_{j}}a_{j}$ in rapidity $s$ ($J = e_{2}e_{3}$ the complex unit of the boundary modes, numbers on the window read in $\mathrm{span}\{1, J\}$, \claimref{synthesis.transverse-bivector-is-the-complex-unit-of-the-radial-atom}), and let $\varrho_{c} = \bigotimes_{j}\operatorname{diag}(1, e^{-2\pi cw_{j}})/(1 + e^{-2\pi cw_{j}})$ in the occupation basis. It is diagonal in the number basis, so it commutes with every number-diagonal generator and with the parity operator (even, stationary, compatible with (L$_E$)); it is the Gibbs state of $2\pi c\sum_{j}w_{j}N_{j}$, hence KMS at rapidity period $2\pi c$ and quasi-free with diagonal symbol. A single mode is KMS at period $2\pi$ only with $n/(1 - n) = e^{-2\pi w}$; the model has $n/(1 - n) = e^{-2\pi cw}$, so for $c \neq 1$ every mode with $w \neq 0$ violates (R$_{\mathrm{loc}}$). The check verifies this with $u = e^{-2\pi}$ as an exact rational-function symbol at weights $1, 2, 3$ and $c = 2$, and that the array weights $-1, 0, 1$ keep the period-$2\pi$ state on a Killing window.

The labels by null direction are used so that the local-horizon clause and its variant are satisfiable at nonzero weight at all: on modes shared by the two sheets of an edge the variant fails for the same reason as the clause (a state KMS at $2\pi c$ for a flow and for its inverse is the trace), and the result filed alongside shows the labels remove that obstruction for every period. With labels, the states of different directions are independent and the product above restricts correctly to every local horizon through the point.

That every other premise holds: the premises (P), (R), (L$_E$), (I), (Q) and the counting premises are statements about the windows of the singular boundary or of a Killing horizon (Theorem \claimref{synthesis.boundary-mode-premises-finite-models} gives their finite readings), and (R$_{\Lambda}$) concerns the cosmological horizon; the model leaves all of these as the store has them. Of the store's 51 postulates only (C$_{\mathrm{loc}}$) and sheaf tameness mention local horizons, diamonds, accelerated observers or all points (search of the postulate texts), and tameness is a regularity condition at points of the degeneration locus, which carries no state. (Loc) fixes the entropy differences $\delta S = \eta\,\delta A$ and says nothing about the state's period: version 2 of the premise does not tie the entropy to the state of the modes (the mode-count reading is excluded there), which is why it holds in the model; under that excluded reading the entropy density would change with $c$. The structural premise of indifference (scrambling.indifference-structural-premise-series) forbids assignments that nothing in the framework determines, but it is applied at the declared level and flagged where it is used; the model's change of period between Killing patches and $U$ is such an assignment, so the premise holds in its declared, scoped reading and not under an unrestricted one. The identification of horizon and cutoff modes concerns the modes of a hole and is kept. So the model satisfies every present premise except (R$_{\mathrm{loc}}$), which it violates for $c \ne 1$; this is independence from the present premises. It needs $U \ne \emptyset$: in exact de Sitter space every local Rindler horizon is at leading order a patch of some observer's cosmological horizon (null hyperplanes of the embedding, edges moved by null rotations), so if (R$_{\Lambda}$) is read for every observer no such $U$ exists there; generic spacetimes and the exterior or interior of a hole away from $r = 2M$ supply one. A premise fixing the state of local-horizon modes, such as the positive-null-energy premise of the result filed alongside (synthesis.positive-null-energy-gives-r-premises), would make the model fail in $U$ and imply (R$_{\mathrm{loc}}$). Imported field-theory theorems that would give period $2\pi$ (Bisognano-Wichmann, Unruh, in Fact \claimref{angle.imported-theorems-deliver-their-hypotheses}) apply to a field theory meeting their hypotheses; they are not premises about the store's modes, and the model is not such a theory.

\emph{(ii)} The geometric half of Theorem \claimref{synthesis.local-boost-temperature-is-local-regularity} is unchanged: the boost field $\chi$ has surface gravity $\kappa$ and rapidity $\kappa$ times its parameter, so a state of rapidity period $2\pi c$ has temperature $\kappa/(2\pi c)$ in the time of $\chi$. On the past half of the generator through $p$, with $\chi = -\kappa\lambda k$ (hypothesis (H2) of Theorem \claimref{synthesis.einstein-equations-from-c-loc-and-clausius}) and $\theta = -\lambda\,\mathrm{Ric}^{\circ}(k,k)|_{p} + O(\lambda^{2})$ (step (2) there),
\[
\delta Q = \int_{\lambda_{0}}^{0}T(k,k)\,(-\kappa\lambda)\,d\lambda = \frac{\kappa}{2}T(k,k)\,\lambda_{0}^{2}, \qquad
\delta A = \int_{\lambda_{0}}^{0}\theta\,d\lambda = \frac12\mathrm{Ric}^{\circ}(k,k)\,\lambda_{0}^{2},
\]
per unit cross-section area, with the error terms of that theorem. So $\delta Q - T_{\chi}\eta\,\delta A = \tfrac{\kappa}{2}\bigl(T(k,k) - \tfrac{\eta}{2\pi c}\mathrm{Ric}^{\circ}(k,k)\bigr)\lambda_{0}^{2} + O(\lambda_{0}^{3})$. The check integrates the two polynomials exactly at three rational parameter sets, carrying $2\pi$ as an independent rational symbol so that no step uses its value, and confirms that the defect vanishes exactly when $T(k,k) = (\eta/2\pi c)\,\mathrm{Ric}^{\circ}(k,k)$.

\emph{(iii)} (H1) at leading order requires $\mathrm{Ric}^{\circ}(k,k) = (2\pi c/\eta)\,T(k,k)$ for every null $k$ at every $p \in U$. The symmetric tensor $S = \mathrm{Ric}^{\circ} - (2\pi c/\eta)T$ then vanishes on all null vectors, so $S = f g$ for a function $f$: the map $S \mapsto (S(k,k))_{k}$ on the ten-dimensional space of symmetric tensors has the one-dimensional kernel spanned by $g$, which the check confirms exactly (rank $9$ over twelve rational null vectors, $g$ in the kernel). From here the argument of Theorem \claimref{synthesis.einstein-equations-from-c-loc-and-clausius} applies word for word with $\eta$ replaced by $\eta/c$: the contracted Bianchi identity and $\nabla^{a}T_{ab} = 0$ then give Einstein's equations with $\Lambda$ an integration constant and $8\pi G_{\mathrm{cl}} = 2\pi c/\eta$, that is $G_{\mathrm{cl}} = c/(4\eta)$. On a Killing horizon the generator $\ell$ has $\theta = \sigma = 0$ along it, so $\mathrm{Ric}^{\circ}(\ell,\ell) = 0$ by the Raychaudhuri equation, and the equation at $k = \ell$ forces $T(\ell,\ell) = 0$ at any period; the null-null equation there is $0 = 0$ whatever $c$ is assigned, and the equations for $k \neq \ell$ determine $S$ by continuity in $k$, so a period that jumps between the Killing patch and the others obstructs nothing.

\emph{(iv)} Theorem \claimref{synthesis.c-loc-entropy-clause-curvature-independent-density}(ii) fixes $\eta = \tfrac14$ from the Schwarzschild family through Theorem \claimref{synthesis.quarter-from-temperature}, which uses the Hawking temperature of (R) at Killing horizons; the model keeps (R), so the calibration is unchanged and $G_{\mathrm{cl}} = c/(4 \cdot \tfrac14) = c$ in units of the $G$ that defines the Planck cell. The check reads $G_{\mathrm{cl}}$ off the defect: it computes the coefficient of $\mathrm{Ric}^{\circ}(k,k)$ in the exactly integrated defect, sets the defect to zero to get $\mathrm{Ric}^{\circ}(k,k) = \mu\,T(k,k)$, and finds $\mu = 8\pi G_{\mathrm{cl}}$ with $G_{\mathrm{cl}} = c$ at $\eta = \tfrac14$ for $c = 1, 2, 3$ (with $2\pi$ carried as a rational symbol). Since only $\eta/c$ enters (ii) and (iii), the field equations do not separate the two clauses: any change of the local period can be compensated by a change of the density, and the Killing calibration is what separates them.
